% ============================================================
% SKELETON -- Part II, Section 6: Response Functions and Stability
% Each \paragraph states the goal and the boxed target result.
% "% TODO" lists the derivation still to write.
% "% FROM" names the result of Sections 1-5 that a step reuses.
% ============================================================

\section{Response Functions and Stability}

\paragraph{Scope of this section.}
Sections~3--5 derive the laws of thermodynamics, the potentials and the
fluctuation formulas from entropy maximization. Experiments measure how one
variable responds to a change of another, not the potentials themselves.
This section defines these responses, relates them through exact identities
and collects the conditions that stability imposes on them.

\subsection{Response Functions}

\paragraph{Measurable second derivatives.}
Each \textcolor{blue}{response function} is a second derivative of a
potential. The heat capacities measure the response of the entropy to the
temperature, the thermal expansion coefficient measures the response of the
volume to the temperature, and the compressibilities measure the response
of the volume to the pressure.
% FROM: C_V (Section 3); delta Q = dH at fixed P (Section 4); G(T, P, N) (Section 5).
% TODO: define C_P = T (dS/dT)_P = (dH/dT)_P, alpha_P, kappa_T and
%       kappa_S = -(1/V)(dV/dP)_S. The subscript P in alpha_P avoids a clash
%       with the phase label alpha of Part III.
\begin{equation}
    \boxed{C_P=-T\frac{\partial^2G}{\partial T^2},\qquad
    \alpha_P=\frac{1}{V}\frac{\partial^2G}{\partial T\,\partial P},\qquad
    \kappa_T=-\frac{1}{V}\frac{\partial^2G}{\partial P^2}.}
\end{equation}

\paragraph{Maxwell relations.}
The mixed second derivatives of a potential do not depend on the order of
differentiation. Each potential therefore gives one
\textcolor{blue}{Maxwell relation}.
% TODO: derive from dU, dF, dH and dG; show how a relation turns an
%       unmeasurable entropy derivative into a measurable one.
\begin{equation}
    \boxed{
    \begin{aligned}
        \left(\frac{\partial T}{\partial V}\right)_{\!S,N}&=-\left(\frac{\partial P}{\partial S}\right)_{\!V,N}, &
        \left(\frac{\partial S}{\partial V}\right)_{\!T,N}&=\left(\frac{\partial P}{\partial T}\right)_{\!V,N},\\
        \left(\frac{\partial T}{\partial P}\right)_{\!S,N}&=\left(\frac{\partial V}{\partial S}\right)_{\!P,N}, &
        \left(\frac{\partial S}{\partial P}\right)_{\!T,N}&=-\left(\frac{\partial V}{\partial T}\right)_{\!P,N}.
    \end{aligned}}
\end{equation}

\paragraph{The cyclic relation.}
Three variables linked by one equation of state satisfy the cyclic
relation. The slope of an isochore in the $P$--$T$ plane then follows from
$\alpha_P$ and $\kappa_T$.
% FROM: the implicit-function result dU/dX|_S = -(dS/dX)/(dS/dU) (Section 4).
\begin{equation}
    \boxed{\left(\frac{\partial P}{\partial T}\right)_{\!V}
    \left(\frac{\partial T}{\partial V}\right)_{\!P}
    \left(\frac{\partial V}{\partial P}\right)_{\!T}=-1,
    \qquad
    \left(\frac{\partial P}{\partial T}\right)_{\!V}=\frac{\alpha_P}{\kappa_T}.}
\end{equation}

\paragraph{Relations between the response functions.}
The response functions are not independent. Two exact relations connect
the heat capacities and the compressibilities.
% TODO: C_P - C_V from S(T, V(T, P)) and the second Maxwell relation;
%       C_P / C_V = kappa_T / kappa_S from the cyclic relation.
\begin{equation}
    \boxed{C_P-C_V=\frac{TV\alpha_P^2}{\kappa_T},\qquad
    \frac{C_P}{C_V}=\frac{\kappa_T}{\kappa_S}.}
\end{equation}

\subsection{Response Functions as Fluctuations}

\paragraph{Heat capacity and energy fluctuation.}
Section~3 relates $C_V$ to the energy fluctuation of the canonical
ensemble. The relation is recalled here as the first member of a family.
% FROM: Var(U) = k_B T^2 C_V (Section 3).
\begin{equation}
    \boxed{\mathrm{Var}(U)=k_BT^2C_V.}
\end{equation}

\paragraph{Compressibility and particle-number fluctuation.}
The fluctuation formula of the grand canonical ensemble and the Gibbs--Duhem
relation give $\kappa_T$ the same statistical meaning.
\textcolor{red}{Here, $n\equiv N/V$ denotes the number density.}
% FROM: Var(N) = k_B T (dN/dmu)_{T,V} and Gibbs--Duhem, N dmu = V dP at
%       fixed T (Section 5).
\begin{equation}
    \boxed{\frac{\mathrm{Var}(N)}{N}=n\,k_BT\,\kappa_T.}
\end{equation}

\subsection{Thermodynamic Stability Revisited}

\paragraph{What Sections 3 and 4 already establish.}
Section~3 shows that the concavity of $S(U)$ is equivalent to $C_V>0$.
Section~4 shows that $U(S,V,N)$ is convex, so $F$ is convex in $V$ and $G$
is concave in $P$. This subsection translates these statements into signs
of the response functions and adds the conditions that Sections~3 and 4 do
not state.
% FROM: d^2S/dU^2 = -1/(C_V T^2) (Section 3); convexity of partial
%       transforms (Section 4).

\paragraph{Positive heat capacities and compressibilities.}
The convexity of $F$ in $V$ gives $\partial^2F/\partial V^2=1/(V\kappa_T)>0$,
and the fluctuation formula gives $\kappa_T>0$ independently. The relations
between the response functions then order the two heat capacities and the
two compressibilities.
% TODO: both routes to kappa_T > 0; C_P >= C_V and kappa_T >= kappa_S.
\begin{equation}
    \boxed{C_P\geq C_V>0,\qquad\kappa_T\geq\kappa_S>0.}
\end{equation}

\paragraph{The principle of Le Chatelier.}
A stable system responds to an imposed change in the direction that reduces
the effect of the change. The positive response functions state this
principle quantitatively.
% TODO: heating at fixed V raises T by dU/C_V; compression at fixed T
%       raises P by dV/(V kappa_T).
